\section{ObserverBench design}

An observer in ObserverBench connects a measurement procedure to a downstream
use. For a task state $h$, we write the estimator and direction provider as
\[
E(h)=\widehat z,
\qquad
D(h)=d.
\]
The controller receives the target and estimate and returns an action magnitude,
\[
u=C(y^\star,\widehat z).
\]
The task owns the plant or intervention rule, initial conditions, measurement
design, and validation metric. This division makes the experimental factors
explicit. A fixed-direction comparison changes only $E$. A fixed-estimator
comparison changes only $D$. A factorial study changes both and reports both
main effects and their interaction.

The earlier diagonal comparison changed $E$ and $D$ together and called the
pair an observer. That composition remains useful operationally, but it cannot
support a causal statement about estimation alone. We therefore use
``estimator--direction pair'' for diagonal arms and reserve ``estimator effect''
for a contrast with direction held fixed.

\paragraph{Fixed task and controller.}
For a benchmark-comparable run, the task fixes the plant, controller settings,
initial-state distribution, intervention operator, admissible direction space,
and normalization. Direction is an experimental factor; it is not silently
treated as part of a fixed actuator. Custom controllers are permitted by the
lower-level composition interface, but those runs are marked non-comparable to
the bundled benchmark arm.

\paragraph{Known or borrowed validation targets.}
Synthetic tasks are useful when they isolate a mechanism. Real-model tasks are
useful when the target is specified before looking at the result. Ctl-2 uses the
trained model's target head as its plant readout. IOI uses published head groups
as borrowed structure. The fitted nuisance probe in Ctl-2 is reported as a
probe displacement, not as independent ground truth.

\paragraph{Held-out validation.}
A diagnostic intervention can show that an interaction exists. It does not show
that an observer predicts new interventions. IOI therefore moves from
whole-group diagnostics to held-out head subsets. Ctl-2 calibrates finite
responses on a separate split before evaluating the closed loop. Holding data
out prevents the recovery method from being judged on the measurements used to
fit it.

\paragraph{Continuous errors before threshold flags.}
Thresholded divergence rates are easy to read, but a trajectory can cross or
miss a threshold because of a small numerical change. Ctl-2 therefore uses
integrated squared error and final target error as primary outcomes. Threshold
rates remain secondary descriptions of the clipped rollout.

\paragraph{Report conditioning and support.}
Identifiability is not enough when the required calibration gain or projected
direction is very large. Every Ctl-2 arm reports the target gain, observer
self-gain, local pole, direction norm, projection scale, clipping, displacement
from the initial state, nearest-clean distance, affine-span distance, and
convex-hull extrapolation. Membership in a clean-state affine span is not called
on-manifold behavior.

\paragraph{Report nulls.}
A workbench must be allowed to recommend the cheap observer. The $\gamma=0$
control null and the IOI random-subset result test that rule directly.
